\setcounter{section}{6}

\Needspace{5\baselineskip}
\hypertarget{ux7a00ux758fux6ce8ux610fux529bsparse-attention}{%
\section{稀疏注意力（Sparse Attention）}\label{ux7a00ux758fux6ce8ux610fux529bsparse-attention}}

\Needspace{5\baselineskip}
\hypertarget{ux4e00ux6838ux5fc3ux601dux60f3-8}{%
\subsection{一、核心思想}\label{ux4e00ux6838ux5fc3ux601dux60f3-8}}

稀疏注意力的核心思想是：并非所有 token 对之间都需要进行交互。通过引入结构性稀疏模式，只计算注意力矩阵中一部分重要元素，可以把全注意力的 \(O(n^2)\) 复杂度降低到 \(O(n\sqrt{n})\)、\(O(n\log n)\) 或近似线性的级别。

\Needspace{5\baselineskip}
\hypertarget{ux4e8cux5178ux578bux4ee3ux8868}{%
\subsection{二、典型代表}\label{ux4e8cux5178ux578bux4ee3ux8868}}

Blockwise Attention：将序列分块，只在块内或特定的块之间计算注意力。

Sliding Window Attention（滑动窗口注意力）：每个 token 只关注其前后一定窗口大小内的邻居 token。这在类似 BERT 这样的编码器中很常见。

Dilated Attention：类似于空洞卷积，在滑动窗口的基础上进行跳跃，以扩大感受野。

Global + Local Attention：指定少数 token（如 \texttt{\textless{}s\textgreater{}}）拥有全局注意力，可以关注所有 token，而其他 token 只进行局部注意力。

应用：Longformer、BigBird 等模型就采用了这种思想，能够处理极长文档。

\Needspace{5\baselineskip}
\hypertarget{ux4e09dsadeepseek-sparse-attentiondeepseek-v3.22025}{%
\subsection{三、DSA（DeepSeek Sparse Attention，DeepSeek-V3.2，2025）}\label{ux4e09dsadeepseek-sparse-attentiondeepseek-v3.22025}}

\Needspace{5\baselineskip}
\hypertarget{ux6838ux5fc3ux601dux60f3ux5148ux7c97ux7b5bux540eux7ec6ux7b97}{%
\subsubsection{1. 核心思想：``先粗筛，后细算''}\label{ux6838ux5fc3ux601dux60f3ux5148ux7c97ux7b5bux540eux7ec6ux7b97}}

DSA 的目标是在不计算所有 Query-Key 点积的情况下，快速找到高价值的 Key。

\Needspace{5\baselineskip}
输入包括：

\begin{itemize}
\tightlist
\item
  Query 侧索引信号：从 \(c_t^Q\) 衍生出的 \(q^I_{t,j}\)。
\item
  Key 侧索引信号：\(k_t^I\)。
\item
  权重或偏置项：\(w^I_{t,j}\)。
\end{itemize}

\Needspace{5\baselineskip}
处理流程包括：

\begin{enumerate}
\def\labelenumi{\arabic{enumi}.}
\tightlist
\item
  部分 RoPE（Partially apply RoPE）：为了降低计算量，索引器不需要全精度的位置编码，只对索引向量应用部分 RoPE。
\item
  Lightning Indexer（闪电索引器）：这是一个轻量级计算模块。它接收粗粒度的 Query 和 Key 索引向量，快速计算粗略的相关性分数。
\item
  Top-k Selector（Top-k 选择器）：基于索引器计算出的分数，动态选出当前 Query 最关注的 \(k\) 个 Key-Value 块。
\end{enumerate}

\Needspace{5\baselineskip}
\hypertarget{ux63a8ux7406ux5de5ux4f5cux6d41}{%
\subsubsection{2. 推理工作流}\label{ux63a8ux7406ux5de5ux4f5cux6d41}}

\Needspace{8\baselineskip}
Step 1：投影降维（Down-Projection）。

\Needspace{5\baselineskip}
模型不直接使用高维的 \(h_t\) 或完整的 \(c_t^Q\) 进行检索，而是通过一个轻量级线性层 \(W^I\) 将输入投影到一个低维空间：

\[
q^I_{t,j} = W^I_Q \cdot c^Q_t
\]

\[
k^I_t = W^I_K \cdot c^{KV}_t
\]

其中，\(q^I_{t,j}\) 是 Query 侧索引向量，\(k^I_t\) 是 Key 侧索引向量。这里的索引向量维度 \(d_{index}\) 远小于正常的 head 维度 \(d_{head}\)。例如，\(d_{head}\) 可能是 128，而 \(d_{index}\) 可能只有 32 或 16，这会显著减少索引阶段的计算量。

\Needspace{8\baselineskip}
Step 2：块级划分（Block-wise Segmentation）。

这里的``粗粒度''不仅指向量维度低，也指时间步的粒度。DSA 通常不会对每一个历史 token 逐个计算分数，而是将 KV Cache 划分为多个固定大小的 Block（块），例如每 64 个 token 为一块。

索引器会为每一个 Block 计算一个代表性的 Key 索引向量，通常可以取该 Block 内所有 \(k^I\) 的均值，也可以对特定位置采样。

\Needspace{8\baselineskip}
Step 3：快速打分（Lightweight Scoring）。

\Needspace{5\baselineskip}
使用点积计算 Query 索引向量与 Block 索引向量之间的相关性分数：

\[
S_{block} = (q^I_{t,j})^T \cdot k^I_{block}
\]

图中提到的 ``partially apply RoPE'' 指的是：为了进一步节省算力，同时保留位置信息，索引器只对向量的一小部分维度应用旋转位置编码，或者使用简化版的位置编码。

\Needspace{8\baselineskip}
Step 4：Top-k 筛选（Gating）。

根据 \(S_{block}\) 的大小，选出分数最高的 \(k\) 个 Block。后续真正的``重型''注意力计算（Core Attention）只针对这 \(k\) 个被选中的 Block 进行加载和计算，其余大量无关信息直接忽略。

\Needspace{5\baselineskip}
之所以这种方法更快，是因为原始全注意力的计算量大致为：

\[
O(L \cdot d_{head})
\]

\Needspace{5\baselineskip}
而索引器的计算量大致为：

\[
O\left(\frac{L}{B} \cdot d_{index}\right)
\]

其中，\(B\) 是块大小。由于 \(B > 1\) 且 \(d_{index} \ll d_{head}\)，索引器开销通常远小于完整注意力计算。

\Needspace{5\baselineskip}
\hypertarget{ux8badux7ec3ux65b9ux6cd5}{%
\subsubsection{3. 训练方法}\label{ux8badux7ec3ux65b9ux6cd5}}

（1）Lightning Indexer 的训练目标。

由于 Top-K 操作本身是不可导的，也就是无法直接通过 Top-K 选择反向传播主模型的预测误差来更新索引器，DeepSeek 为 Lightning Indexer 设计了独立的监督信号。

Lightning Indexer 被训练来预测每个 token 的 Query 与每个 token 的 Key 之间的注意力权重。训练目标是模仿原始全注意力（Dense Attention）的权重分布，准确说是所有注意力头的权重平均分布。

\Needspace{5\baselineskip}
损失函数可以理解为 \(n\) 个 KL 散度之和。对于第 \(i\) 个 token，它对第 \(j\) 个 token 的注意力分数（\(j = 1, 2, \ldots, n\)）构成概率分布 \(p(j \mid i)\)。再计算索引器分布 \(p_{Indexer}(j \mid i)\) 与稠密注意力分布 \(p_{Dense}(j \mid i)\) 之间的 KL 散度，并对 \(i = 1, 2, \ldots, n\) 累加：

\[
\mathcal{L}_{indexer}
= \sum_{i=1}^{n} D_{KL}\left(p_{Dense}(\cdot \mid i) \,\|\, p_{Indexer}(\cdot \mid i)\right)
\]

（2）整体训练步骤。

第一步：在全注意力（Dense Attention）下，训练主模型。

第二步：冻结主模型，保持全注意力开启，初始化训练 Lightning Indexer，使之与主模型的注意力分布对齐。

第三步：开启稀疏化（Top-K Selection），同步用各自的损失函数训练主模型和 Lightning Indexer。

\FloatBarrier
\Needspace{5\baselineskip}
\hypertarget{ux53c2ux8003ux6587ux732e-97}{%
\subsection{参考文献}\label{ux53c2ux8003ux6587ux732e-97}}

\vspace{0.25em}
\begin{itemize}
\tightlist
\item
  Qiu, J., Ma, H., Levy, O., Yih, S. W., Wang, S., \& Tang, J. (2019). \href{https://arxiv.org/abs/1911.02972}{Blockwise Self-Attention for Long Document Understanding}. arXiv:1911.02972.
\item
  Child, R., Gray, S., Radford, A., \& Sutskever, I. (2019). \href{https://arxiv.org/abs/1904.10509}{Generating Long Sequences with Sparse Transformers}. arXiv:1904.10509.
\item
  Beltagy, I., Peters, M. E., \& Cohan, A. (2020). \href{https://arxiv.org/abs/2004.05150}{Longformer: The Long-Document Transformer}. arXiv:2004.05150.
\item
  Zaheer, M., Guruganesh, G., Dubey, K. A., et al.~(2020). \href{https://arxiv.org/abs/2007.14062}{Big Bird: Transformers for Longer Sequences}. NeurIPS 2020.
\item
  Ding, J., Ma, S., Dong, L., et al.~(2023). \href{https://arxiv.org/abs/2307.02486}{LongNet: Scaling Transformers to 1,000,000,000 Tokens}. arXiv:2307.02486.
\item
  DeepSeek-AI. (2025). \href{https://arxiv.org/abs/2512.02556}{DeepSeek-V3.2: Pushing the Frontier of Open Large Language Models}. arXiv:2512.02556.
\end{itemize}

\clearpage
